\documentclass[12pt,aspectratio=169,english,finnish,pdflatex%,handout
]{beamer}
\definecolor{MyGreen}{RGB}{50, 120, 50}
\usecolortheme[named=MyGreen]{structure}

\usepackage{babel}
\usepackage[utf8]{inputenc}
\usepackage[T1]{fontenc}
\usepackage{amsmath,amssymb} 
\usepackage{animate}
\usepackage{multimedia}

\usepackage{natbib}
\bibpunct[: ]{(}{)}{,}{}{}{;}

\usepackage{tikz}

\usepackage{tipa}

\usepackage{hyperref}

\setbeamertemplate{navigation symbols}{}

\graphicspath{{figures/}}

\setlength{\leftmargini}{0pt}
\setlength{\leftmarginii}{1em}

%% Write out the names of graphics files being included 
\newwrite\graphics
\immediate\openout\graphics=\jobname.graphics%
\let\oincludegraphics\includegraphics% store original \includegraphics
\renewcommand{\includegraphics}[2][]{% prepend to it (could also use xpatch, etc.)
  \immediate\write\graphics{#2}
  \oincludegraphics[#1]{#2}}

\newcommand\Wider[2][3em]{%
\makebox[\linewidth][c]{%
  \begin{minipage}{\dimexpr\textwidth+#1\relax}
  \raggedright#2
  \end{minipage}%
  }%
}

\newcommand{\kommentti}[1]{
  {\bf[#1]}
}


\title{CSD 598 - Winter 2026
  \\~\\
  An example of exploration and inference
}
\author{Pertti Palo} 
\date{24 Feb 2026} 


\begin{document}

\frame{\titlepage
  \centering
} 


\frame{\frametitle{First reflection I: Ethics}
  Any questions? Any thoughts? What did you find most difficult about this
  assignment?
 }


\frame{\frametitle{First reflection II: Ethics}
  \begin{itemize}
    \item Risk came up a lot unsurprisingly.
    \item I was impressed with the level of careful consideration you showed of
    different potential risk factors.
    \item You can also question the training itself. One of you brought up the
    absence of AI in the CORE materials.
    \item Contextualisation for the reader is good to do, but it's not
    absolutely necessary if the story of thinking is clear without the details.
    'Keep the grader happy' is a good general rule in this context.
  \end{itemize}
}


\frame{\frametitle{Last time: Visualisation to inferential statistics}
  Any questions? Any thoughts?
}


\frame{\frametitle{Exploration and inference example}
  \begin{itemize}
    \item This comes from the latest data analysis I have run.
    \item As it is unpublished, I ask it not be recorded, and hence will not
    describe it here on the slides.
    \item I am happy to talk about it in detail and show how I've done the
    whole thing.
  \end{itemize}
}


\frame{\frametitle{Model is a test, test is a model -- sometimes}
  \begin{itemize}
    \item ANOVA can be seen as a single predictor (controlled variable) linear
      regression model where the only predictor is a categorical variable. 
    \item Simple regression -- line fitting -- can be seen as another way of
    testing for linear (Pearson's) correlation.
  \end{itemize}
}


\frame{\frametitle{Model diagnosis -- this will likely be moved to next time}
  \begin{itemize}
    \item Collinearity (linear correlation) of the independent/predicting
    variables is bad (mathematically). 
    \begin{itemize}
      \item It can usually be spotted fairly easily by calculating correlations
        among the predictors and/or making scatter plots like a pairs plot.
    \end{itemize}
    \item Heteroscedasticity will mess more the efficiency of prediction than
    the accuracy of it.
    \item Outliers can mess prediction accuracy.
    \item Here a bit of data exploration and model diagnostic plots go a long
    way:
    \begin{itemize}
      \item Neat explanation of diagnostic plots can be found here
      \url{https://library.virginia.edu/data/articles/diagnostic-plots}
      \item Collinearity can be seen in for example a pairs plot or a
      correlation matrix (not to be confused with e.g. "Correlation of Fixed
      Effects" in lmer output).
    \end{itemize}
  \end{itemize}
}


\frame{\frametitle{Couple of nice resources}
  \begin{itemize}
    \item A detailed look on ANOVA \url{https://statsandr.com/blog/anova-in-r/}. 
    \item And ANOVA's cousin Kruskal-Wallis \url{https://statsandr.com/blog/kruskal-wallis-test-nonparametric-version-anova/}
  \end{itemize}
}


% \frame{\frametitle{References}
  
% \bibliographystyle{apalike}
% \bibliography{science_combined}

% }
\end{document}

